\documentclass[../../main.tex]{subfiles}

\begin{document}

\section{可数集与不可数集}\label{section:ch01-可数不可数}

在预备知识中我们称两个集合 $A$ 和 $B$ 对等, 若存在将 $A$ 映上 $B$ 的一对一映射 $f$. 我们称这样的 $f$ 为集合 $A$ 和 $B$ 之间的一一对应. 对等在集合间定义了一个等价关系, 即它是自反的、对称的与传递的（见习题 20）. 记自然数 $\{k\in\mathbf{N}\mid 1\leqslant k\leqslant n\}$ 为 $\{1,\cdots,n\}$ 是方便的. 关于对等的第一个观察是, 对任何自然数 $n$ 和 $m$, 集合 $\{1,\cdots,n+m\}$ 与集合 $\{1,\cdots,n\}$ 不对等. 该观察常称为鸽笼原理, 可用关于 $n$ 的归纳法证明（见习题 21）.

\begin{definition}[有限、可数无穷、可数集]\label{definition:ch01-可数}
集合 $E$ 称为是有限的, 若它或者是空集, 或者存在自然数 $n$ 使得 $E$ 与 $\{1,\cdots,n\}$ 对等. 我们说 $E$ 是可数无穷的, 若 $E$ 与自然数集 $\mathbf{N}$ 对等. 有限或可数无穷的集合称为可数集. 不是可数的集合称为不可数集.
\end{definition}

观察到若一个集与可数集对等, 则它是可数的. 在下面定理的证明中我们将用鸽笼原理和\cref{theorem:ch01-自然数最小元}, 它告诉我们每个非空自然数集有一个最小元或首元.

\begin{theorem}\label{theorem:ch01-可数子集}
可数集的子集是可数的. 特别是, 每个自然数集是可数的.
\end{theorem}

\begin{proof}
令 $B$ 为可数集而 $A$ 是 $B$ 的一个非空子集. 首先考虑 $B$ 是有限的情形. 令 $f$ 为 $\{1,\cdots,n\}$ 与 $B$ 之间的一一对应. 定义 $g(1)$ 为第一个使得 $f(j)$ 属于 $A$ 的自然数 $j$, $1\leqslant j\leqslant n$. 由于 $f\circ g$ 是 $\{1\}$ 与 $A$ 之间的一一对应, 若 $A=\{f(g(1))\}$, 证明完成. 否则, 定义 $g(2)$ 为使得 $f(j)$ 属于 $A\setminus\{f(g(1))\}$ 的第一个自然数 $j$, $1\leqslant j\leqslant n$. 鸽笼原理告诉我们至多 $N$ 步后该归纳选择过程终止, 其中 $N\leqslant n$. 因此 $f\circ g$ 是 $\{1,\cdots,N\}$ 与 $A$ 之间的一一对应. 于是 $A$ 有限.

现在考虑 $B$ 是可数无穷的情形. 令 $f$ 为 $\mathbf{N}$ 与 $B$ 之间的一一对应. 定义 $g(1)$ 为第一个使得 $f(j)$ 属于 $A$ 的自然数 $j$. 如同第一种情形的证明, 我们看到若该选择过程终止, 则 $A$ 是有限的. 否则, 该选择过程不终止而 $g$ 在所有的 $\mathbf{N}$ 上恰当定义. 显然 $f\circ g$ 是一一映射, 其中定义域是 $\mathbf{N}$ 而象包含于 $A$ 中. 归纳论证表明对所有 $j$, $g(j)\geqslant j$. 对每个 $x\in A$, 存在某个 $k$ 使得 $x=f(k)$. 因此 $x$ 属于集合 $\{f(g(1)),\cdots,f(g(k))\}$. 因此 $f\circ g$ 的象是 $A$. 因此 $A$ 是可数无穷.
\end{proof}

\begin{corollary}\label{corollary:ch01-可数推论}
下面的集合是可数无穷的:
\begin{enumerate}[(i)]
\item 对每个自然数 $n$, 笛卡儿积 $\mathbf{N}\times\cdots\times\mathbf{N}$（$n$ 次）.
\item 有理数集 $\mathbf{Q}$.
\end{enumerate}
\end{corollary}

\begin{proof}
我们对 $n=2$ 证明 (i), 而一般情形留作归纳法的练习. 定义从 $\mathbf{N}\times\mathbf{N}$ 到 $\mathbf{N}$ 的映射 $g$ 为 $g(m,n)=(m+n)^{2}+n$. 映射 $g$ 是一对一的. 事实上, 若 $g(m,n)=g(m',n')$, 则 $(m+n)^{2}-(m'+n')^{2}=n'-n$, 因此
\begin{align*}
|m+n+m'+n'|\cdot|m+n-m'-n'|=|n'-n|.
\end{align*}
若 $n\neq n'$, 则自然数 $m+n+m'+n'$ 大于自然数 $|n'-n|$, 这是不可能的. 于是 $n=n'$, 因而 $m=m'$. 因此 $\mathbf{N}\times\mathbf{N}$ 与可数集 $\mathbf{N}$ 的子集 $g(\mathbf{N}\times\mathbf{N})$ 对等. 我们从前一个定理推出 $\mathbf{N}\times\mathbf{N}$ 是可数的. 为证明 $\mathbf{Q}$ 的可数性, 我们首先从素数分解定理推出每个正有理数 $x$ 可唯一写成 $x=p/q$, 其中 $p$ 和 $q$ 是互素的自然数. 对 $x=p/q>0$ 定义从 $\mathbf{Q}$ 到 $\mathbf{N}$ 的映射 $g$ 为 $g(x)=2((p+q)^{2}+q)$, 其中 $p$ 和 $q$ 是互素的自然数, $g(0)=1$, 而对 $x<0$, $g(x)=g(-x)+1$. 我们将证明 $g$ 是一对一的留作练习. 于是 $\mathbf{Q}$ 与 $\mathbf{N}$ 的一个子集对等, 因此根据前一个定理, 是可数的. 我们将用鸽笼原理证明 $\mathbf{N}\times\mathbf{N}$ 和 $\mathbf{Q}$ 都不是有限的留作练习.
\end{proof}

对于可数无穷集 $X$, 我们说 $\{x_{n}\mid n\in\mathbf{N}\}$ 是 $X$ 的一个列举, 若
\begin{align*}
X=\{x_{n}\mid n\in\mathbf{N}\},\qquad x_{n}\neq x_{m}\text{（若 }n\neq m\text{）}.
\end{align*}

\begin{theorem}\label{theorem:ch01-可数象}
非空集是可数的当且仅当它是某个定义域为非空可数集的函数的象.
\end{theorem}

\begin{proof}
令 $A$ 为非空可数集, 而 $f$ 为将 $A$ 映上 $B$ 的映射. 假定 $A$ 是可数无穷的, 而将有限的情形留作练习. 通过 $A$ 与 $\mathbf{N}$ 之间的一一对应的复合, 我们可以假定 $A=\mathbf{N}$. 定义 $A$ 中的两点 $x,x'$ 为等价的, 若 $f(x)=f(x')$. 这是一个等价关系, 即它是自反的、对称的与传递的. 令 $E$ 为 $A$ 的子集, 它由每个等价类的一个成员组成. 则 $f$ 在 $E$ 的限制是 $E$ 与 $B$ 之间的一一对应. 但 $E$ 是 $\mathbf{N}$ 的子集, 因此, 根据\cref{theorem:ch01-可数子集}, 是可数的. 集合 $B$ 与 $E$ 对等, 因此 $B$ 是可数的. 逆断言是显然的, 若 $B$ 是非空可数集, 则它或者与自然数的一个初始部分对等, 或者与自然数全体对等.
\end{proof}

\begin{corollary}\label{corollary:ch01-可数族并}
可数集的可数族的并是可数的.
\end{corollary}

\begin{proof}
令 $\Lambda$ 为可数集且对每个 $\lambda\in\Lambda$, 令 $E_{\lambda}$ 为可数集. 我们将证明 $E=\bigcup_{\lambda\in\Lambda}E_{\lambda}$ 是可数的. 若 $E$ 是空集, 则它是可数的. 因此我们假设 $E\neq\varnothing$. 我们考虑 $\Lambda$ 是可数无穷的情形, 而将有限的情形留作练习. 令 $\{\lambda_{n}\mid n\in\mathbf{N}\}$ 为 $\Lambda$ 的一个列举. 固定 $n\in\mathbf{N}$. 若 $E_{\lambda_{n}}$ 是有限且非空的, 选取自然数 $N(n)$ 与将 $\{1,\cdots,N(n)\}$ 映上 $E_{\lambda_{n}}$ 的一一映射 $f_{n}$; 若 $E_{\lambda_{n}}$ 是可数无穷的, 选取 $\mathbf{N}$ 映上 $E_{\lambda_{n}}$ 的一一映射 $f_{n}$. 定义
\begin{align*}
E'=\{(n,k)\in\mathbf{N}\times\mathbf{N}\mid E_{\lambda_{n}}\text{ 是非空的, 且若 }E_{\lambda_{n}}\text{ 也是有限的, }1\leqslant k\leqslant N(n)\}.
\end{align*}
定义 $E'$ 到 $E$ 的映射 $f$ 为 $f(n,k)=f_{n}(k)$. 则 $f$ 是 $E'$ 映上 $E$ 的映射. 然而, $E'$ 是可数集 $\mathbf{N}\times\mathbf{N}$ 的子集, 因此, 根据\cref{theorem:ch01-可数子集}, 是可数的.\cref{theorem:ch01-可数象}告诉我们 $E$ 也是可数的.
\end{proof}

我们称实数的区间为退化的, 若它是空的或包含一个单独的成员.

\begin{theorem}\label{theorem:ch01-区间不可数}
一个非退化实数区间是不可数的.
\end{theorem}

\begin{proof}
令 $I$ 为实数的非退化区间. 显然 $I$ 不是有限的. 我们用反证法证明 $I$ 是不可数的. 假定 $I$ 是可数无穷的. 令 $\{x_{n}\mid n\in\mathbf{N}\}$ 为 $I$ 的一个列举. 令 $[a_{1},b_{1}]$ 为 $I$ 的不包含 $x_{1}$ 的非退化的闭有界子区间. 接着令 $[a_{2},b_{2}]$ 为 $[a_{1},b_{1}]$ 的非退化的闭有界子区间, 它不包含 $x_{2}$. 我们归纳地选取非退化闭有界区间的可数族 $\{[a_{n},b_{n}]\}_{n=1}^{\infty}$, 对每个 $n$, $[a_{n+1},b_{n+1}]\subseteq[a_{n},b_{n}]$, 并使得对每个 $n$, $x_{n}\notin[a_{n},b_{n}]$. 非空集 $E=\{a_{n}\mid n\in\mathbf{N}\}$ 有上界 $b_{1}$. 完备性公理告诉我们 $E$ 有上确界. 定义 $x^{*}=\sup E$. 由于 $x^{*}$ 是 $E$ 的一个上界, 对所有 $n$, $a_{n}\leqslant x^{*}$. 另一方面, 由于 $\{[a_{n},b_{n}]\}_{n=1}^{\infty}$ 是下降的, 对每个 $n$, $b_{n}$ 是 $E$ 的上界. 于是, 对每个 $n$, $x^{*}\leqslant b_{n}$. 因此对每个 $n$, $x^{*}$ 属于 $[a_{n},b_{n}]$. 但 $x^{*}$ 属于 $[a_{1},b_{1}]\subseteq I$, 因此存在自然数 $n_{0}$ 使得 $x^{*}=x_{n_{0}}$. 由于 $x^{*}=x_{n_{0}}$ 不属于 $[a_{n_{0}},b_{n_{0}}]$, 我们得到矛盾. 因此, $I$ 是不可数的.
\end{proof}

\begin{exercise}\label{exercise:ch01-16}
证明整数集合 $\mathbf{Z}$ 是可数的.
\end{exercise}

\begin{exercise}\label{exercise:ch01-17}
证明一个集合 $A$ 是可数的, 当且仅当存在一个 $A$ 到 $\mathbf{N}$ 的一一对应.
\end{exercise}

\begin{exercise}\label{exercise:ch01-18}
用归纳法完成\cref{corollary:ch01-可数推论}中 (i) 的证明.
\end{exercise}

\begin{exercise}\label{exercise:ch01-19}
在可数集的有限簇的情形证明\cref{corollary:ch01-可数族并}.
\end{exercise}

\begin{exercise}\label{exercise:ch01-20}
令 $f:A\to B$ 与 $g:B\to C$ 都是一对一与映上的. 证明复合 $g\circ f:A\to B$ 与逆 $f^{-1}:B\to A$ 也是一对一与映上的.
\end{exercise}

\begin{exercise}\label{exercise:ch01-21}
用归纳法证明鸽笼原理.
\end{exercise}

\begin{exercise}\label{exercise:ch01-22}
证明自然数集的所有子集组成的集簇 $2^{\mathbf{N}}$ 是不可数的.
\end{exercise}

\begin{exercise}\label{exercise:ch01-23}
证明可数集的有限簇的笛卡儿积是可数的. 用前一个习题证明从 $\mathbf{N}$ 到 $\mathbf{N}$ 的所有映射族 $\mathbf{N}^{\mathbf{N}}$ 不是可数的.
\end{exercise}

\begin{exercise}\label{exercise:ch01-24}
证明实数的非退化区间不是有限的.
\end{exercise}

\begin{exercise}\label{exercise:ch01-25}
证明任何两个实数的非退化区间是对等的.
\end{exercise}

\begin{exercise}\label{exercise:ch01-26}
集合 $\mathbf{R}\times\mathbf{R}$ 是否与 $\mathbf{R}$ 对等?
\end{exercise}

\end{document}